\section{Experimental platform and methodology}\label{sec:platform}

The regime coordinates are meaningful only if the cache, memory, and communication hierarchies can be measured directly. Every CPU result in this thesis was collected on one bare-metal workstation (``puffin''), an AMD Ryzen Threadripper 3960X with 24 cores in a single socket and NUMA node. Its 8 \emph{core complexes} (CCX) each contain 3 cores that share one 16\,MiB L3 slice, giving 128\,MiB in total. Each core has a 32\,KiB L1d cache, a 512\,KiB L2 cache, a nominal maximum clock of 3.8\,GHz, and a second hardware thread through simultaneous multithreading (SMT). Main memory is 4-channel DDR4-3200 with a theoretical 102.4\,GB/s roof. Puffin also hosts an NVIDIA RTX~3090 with 24\,GB, one half of the GPU comparison below. Its derived constants are centralized in \texttt{include/machine.hpp} under the key \texttt{amd-3960x}; the regime study computes every predicted CPU coordinate from them.

\paragraph{The GPU platform.} The GPU study of \Cref{sec:gpu} adds a second machine, ``synge,'' a two-node cluster carrying a pair of NVIDIA Tesla V100-PCIE-16GB cards per node (four in total). The two cards within a node communicate through PCIe across a cross-socket hop with no NVLink, and the nodes use a 100\,Gb/s Omni-Path fabric. The NVIDIA Collective Communications Library (NCCL) provides GPU-aware collective and point-to-point operations; on Synge it selects an InfiniBand-verbs transport over that fabric. Synge's V100 and puffin's RTX~3090 form a useful but not controlled architectural contrast. Their measured bandwidths and nominal L2 capacities are similar, while their FP64 peaks differ by about $11\times$; they also differ in SM organization, cache hierarchy, clocks, memory technology, ECC behavior, schedulers, and launch behavior. Conclusions from the pair are therefore kernel- and architecture-specific rather than causal estimates of FP64 throughput alone. Their constants live in a separate header (\texttt{include/gpu\_machine.hpp}), typed distinctly from the CPU constants so that a CPU value reaching a GPU coordinate is a compile error rather than a plausible wrong number, and are calibrated per card on an idle device rather than read from datasheets.

\paragraph{Workstation selection.} A cloud arm was considered because cloud providers offer server-grade hardware at reasonable hourly cost, larger problem configurations, and an environment relevant to commercial research. It was omitted because co-tenancy can produce timing variation that cannot be separated from genuine memory-bound behavior inside the guest. Virtual topology is a larger limitation because the hypervisor reports a socket and private last-level-cache slice per virtual CPU rather than the physical cache hierarchy. The regime coordinates depend on that hierarchy through cache capacity and the boundaries crossed by a reduction tree, so they cannot be evaluated reliably on the reported topology. The study therefore requires an observable physical topology and explicit contention controls. Idle-device gating supplies the corresponding safeguard on the shared cluster. Bare metal removes the hypervisor but not operating-system scheduling noise, so timed quantities are reported as medians with the stated spread over repeated trials. On a boost-enabled processor, the median represents a typical repeat, whereas the minimum can reflect a transient boost state.

\paragraph{Reproducibility.} Each reported result can be traced from the implementation to its figure. The repository contains the CMake build, C++23 experiment programs, centralized machine descriptions, run scripts, and a named plotting script for each figure with inline dependency metadata. Tables and captions record the relevant parameters and timing protocols, and the executables use the same audited constants that define the regime coordinates. Additional run metadata, raw measurements, and profiler outputs are available from the author on request.
